\section{Exact Weil-Block Projective Capacity}
  \label{sec:exact-capacity}

  \subsection{Exact block decomposition}

    The three-fold tensor product $\rho_{\mathrm{aff}}^{\otimes 3}$ decomposes into $q^3$
    irreducible blocks $H_{c_1,c_2,c_3}$, each of dimension $q$ (for generic multi-indices)
    or smaller (for degenerate multi-indices).
    We restrict attention to generic blocks in the sense of Definition~\ref{def:generic}.
    For a generic block $H_c$ with multi-index $c=(c_1,c_2,c_3)$, the group action is
    \begin{equation}
      \rho_c(g_1,g_2,g_3) = \rho_{c_1}(g_1)\otimes\rho_{c_2}(g_2)\otimes\rho_{c_3}(g_3),
      \quad g_i=(a_i,b_i,\gamma_i)\in G_q,
      \label{eq:block-action}
    \end{equation}
    where each factor $\rho_{c_i}$ is the exact Weil representation of~\eqref{eq:weil-exact}.
    The ambient space of $H_c$ is $L^2(\mathbb{Z}/q\mathbb{Z})^{\otimes 3}$, of dimension~$q^3$,
    but the image of the block map~\eqref{eq:block-action} lives in a $q$-dimensional
    irreducible subspace (see Proposition~\ref{prop:block-dim} below).

    \begin{proposition}[Block dimension]
      \label{prop:block-dim}
      For a generic multi-index $c=(c_1,c_2,c_3)$, the irreducible block $H_c$ has dimension~$q$.
    \end{proposition}

    \begin{proof}
      Each $\rho_{c_i}$ is an irreducible $q$-dimensional representation of $G_q$
      (the Weil representation at central character $c_i\neq 0$).
      By the Stone--von Neumann theorem for finite Heisenberg groups, every irreducible
      representation of $G_q$ at a fixed non-trivial central character $c\neq 0$ is
      equivalent to $\rho_c$ and has dimension~$q$.
      The tensor product $\rho_{c_1}\otimes\rho_{c_2}\otimes\rho_{c_3}$ decomposes, for
      generic $(c_1,c_2,c_3)$ with $c_1+c_2+c_3\not\equiv 0\pmod{q}$, as a sum of $q^2$
      copies of the irreducible representation $\rho_{c_1+c_2+c_3}$, each of dimension~$q$.
      Hence each constituent block $H_c$ has dimension~$q$.
    \end{proof}

  \subsection{Incremental block capacity: definition}

    Fix a generic block $H_c$.
    For each BFS shell $S_n = B_n\setminus B_{n-1}$ of $\mathrm{Cay}(G_q,S)$, let
    $F_n^{(c)}\subset H_c$ denote the set of projected fingerprint vectors
    \begin{equation}
      F_n^{(c)}
      = \bigl\{\,\Pi_c\bigl(v(\pi,g)\bigr) : g\in S_n,\;\pi\in S^3\,\bigr\},
    \end{equation}
    where $v(\pi,g)=(g s_1,\,g s_1 s_2,\,g s_1 s_2 s_3)\in G_q^3$ is the endpoint triple of
    path $\pi=(s_1,s_2,s_3)$ starting at~$g$, and $\Pi_c: L^2(G_q^3)\to H_c$ is the
    orthogonal projection onto the block~$H_c$ via~\eqref{eq:block-action}.
    Define the subspaces
    \begin{equation}
      W_{<n}^{(c)} = \mathrm{span}\bigl(F_0^{(c)}\cup\cdots\cup F_{n-1}^{(c)}\bigr),
      \qquad
      W_{\leq n}^{(c)} = \mathrm{span}\bigl(F_0^{(c)}\cup\cdots\cup F_n^{(c)}\bigr).
    \end{equation}
    The \emph{incremental rank contribution} of shell $S_n$ in block $H_c$ is
    \begin{equation}
      \Delta r_n^{(c)} = \dim W_{\leq n}^{(c)} - \dim W_{<n}^{(c)},
    \end{equation}
    and the \emph{exact block-wise projective capacity} is
    \begin{equation}
      \Sigma_n^{(c)} = \frac{\Delta r_n^{(c)}}{|S_n|}.
      \label{eq:exact-cap}
    \end{equation}
    Since $\Delta r_n^{(c)}\leq\min(\dim H_c, |S_n|)\leq\min(q, |S_n|)$, we have
    $\Sigma_n^{(c)}\in[0,1]$ whenever $|S_n|\geq 1$.
    The pre-saturation regime is the range of depths $n$ for which $\Sigma_n^{(c)}$ is still
    positive and decaying: the block $H_c$ is saturated at depth $n^*$ defined by
    $\dim W_{\leq n^*}^{(c)}=q$.

  \subsection{Mean capacity and the decay hypothesis}

    Let $\mathcal{C}_{\mathrm{gen}}$ denote the set of sampled generic multi-indices
    used in the numerical experiment (see Section~\ref{sec:sampling} for the sampling protocol).
    The \emph{mean block capacity} is
    \begin{equation}
      \bar{\Sigma}_n = \frac{1}{|\mathcal{C}_{\mathrm{gen}}|}
      \sum_{c\in\mathcal{C}_{\mathrm{gen}}} \Sigma_n^{(c)}.
    \end{equation}
    In the proxy framework of O11~\cite{Beau2026a15}, the mean capacity over generic channels
    satisfies $\bar{\Sigma}_n\sim C\,n^{-\hat{\delta}_{\mathrm{cap}}}$ with
    $\hat{\delta}_{\mathrm{cap}}\approx 3.4$.
    The exact analogue for the full Weil representation is:

    \begin{hypothesis}[Block decay law]
      \label{hyp:decay}
      There exists an exponent $\delta_{\mathrm{exact}}>0$ and constants $C_c>0$ such that
      for all generic blocks $H_c\in\mathcal{C}_{\mathrm{gen}}$,
      \begin{equation}
        \Sigma_n^{(c)} \sim C_c\,n^{-\delta_{\mathrm{exact}}}
        \qquad (n_0\leq n\leq n_1),
        \label{eq:decay-law}
      \end{equation}
      with a common exponent $\delta_{\mathrm{exact}}$ across all generic blocks
      (but block-dependent amplitudes $C_c$).
    \end{hypothesis}

    The universality of the exponent across blocks is motivated by the fact that the Weil
    blocks diagonalise the affine action of $G_q$: each generic block sees the same
    structural constraints from the BFS exploration, so a common dynamical rate is expected.
    This is however an empirical statement that the present paper tests.

  \subsection{Measurability conditions (E1)--(E3)}

    The following conditions, stricter analogues of (B1)--(B3) from O11~\cite{Beau2026a15},
    must all hold for the extraction of $\hat{\delta}_{\mathrm{exact}}$ to be considered reliable.

    \begin{definition}[Measurability conditions]
      \label{def:conditions}
      Let $[n_0,n_1]\subset[1,n^*]$ be the fitting window.
      \begin{itemize}
        \item[\textbf{(E1)}] \textbf{Window depth.}
        The pre-saturation window contains at least 15 data points: $n_1 - n_0 \geq 14$.
        \item[\textbf{(E2)}] \textbf{Inter-block homogeneity.}
        The inter-block variance of the capacity is small relative to the mean squared capacity
        throughout the fitting window:
        \begin{equation}
          \frac{\mathrm{Var}_{c\in\mathcal{C}_{\mathrm{gen}}}(\Sigma_n^{(c)})}{\bar{\Sigma}_n^2}
          < \varepsilon_E
          \qquad \forall\,n\in[n_0,n_1],
        \end{equation}
        with $\varepsilon_E = 1$.
        If (E2) is satisfied, the decay exponent is universal across blocks in the sense of
        Hypothesis~\ref{hyp:decay}.
        If (E2) is violated, the blocks fragment dynamically and the mean exponent
        $\hat{\delta}_{\mathrm{exact}}$ captures only an average over block-dependent rates.
        \item[\textbf{(E3)}] \textbf{Goodness of fit.}
        The log-log regression of $\bar{\Sigma}_n$ against $n$ over $[n_0,n_1]$ achieves
        $R^2\geq 0.97$.
      \end{itemize}
    \end{definition}

    Conditions (E1)--(E3) are strictly stronger than (B1)--(B3) of O11 in the following sense:
    (E2) now measures variance across exact Weil blocks (not proxy channels), so its
    satisfaction implies universality at the level of the exact irreducible decomposition.
    A violation of (E2) that was not present at the proxy level would indicate that the
    central coordinate~$\gamma$ introduces genuine inter-block heterogeneity.

  \subsection{The Block Independence Hypothesis}
    \label{sec:bih}

    Hypothesis~\ref{hyp:decay} can be reformulated operationally as:

    \begin{hypothesis}[Block Independence Hypothesis (BIH)]
      \label{hyp:bih}
      Under conditions (E1)--(E3), the exact block-wise projective capacities $\Sigma_n^{(c)}$
      converge to a common decay exponent $\delta_{\mathrm{exact}}$ across all generic blocks
      $H_c\in\mathcal{C}_{\mathrm{gen}}$.
      The blocks are therefore said to be \emph{dynamically independent at leading order}:
      their individual amplitudes $C_c$ may differ, but their decay rates coincide.
    \end{hypothesis}

    This is labelled a hypothesis, not a theorem, because its truth depends on the
    outcome of the numerical experiment.
    The Block Independence Hypothesis is the central empirical claim of O12.
    If it holds, condition (E2) is satisfied and a single exponent $\hat{\delta}_{\mathrm{exact}}$
    can be extracted from the mean capacity.
    If it fails, the decay is block-dependent and one must instead report a distribution
    of per-block exponents $\hat{\delta}_c$.
    In either case the result is scientifically informative: confirmation supports the
    universality of the Cosmochrony spectral framework at the exact representation level,
    while failure would identify a new structural obstruction requiring analysis.

    \begin{remark}[Universality in mean vs block-by-block]
      \label{rem:universality}
      There is a subtle distinction between two levels of universality.
      \emph{Block-by-block universality} holds when every individual curve $\Sigma_n^{(c)}$
      decays with the same exponent.
      \emph{Mean universality} holds only when the mean $\bar{\Sigma}_n$ decays as a power
      law, even if individual blocks deviate.
      Condition (E2) tests block-by-block universality.
      If (E2) is violated but $R^2\geq 0.97$ still holds for $\bar{\Sigma}_n$, the result
      is a mean universality statement only, and the per-block dispersion must be reported.
    \end{remark}
